\section{Qupit Clifford Circuits and the Main Theorems}\label{sec:background}

Fix an odd prime $p$ and $\omega = e^{2\pi i/p}$.  On the basis
$\ket{0},\dots,\ket{p-1}$ indexed by $\Zp$, $X\ket{j} = \ket{j+1}$ and
$Z\ket{j} = \omega^{j}\ket{j}$ generate, with the phases $\omega^t$, the
Pauli group $\Pauli$, and $H\colon \ket{j} \mapsto
\tfrac{\delta_p}{\sqrt p}\sum_k \omega^{jk}\ket{k}$, $S\colon \ket{j}
\mapsto \omega^{j(j-1)/2}\ket{j}$, $CZ\colon \ket{j}\ket{l} \mapsto
\omega^{jl}\ket{j}\ket{l}$ generate, with the phases $\omega^t$, the
Clifford group $\Cliff$ ($\delta_p = i^{(p-1)/2}$ makes the
determinant of $H$ equal to $1$; the determinants of $S$ and $CZ$ are
$1$ too)~\cite{bian2026qudit}.  Modulo scalars, $\PCliff =
\Cliff/\langle\omega\rangle$ and $\PPauli \cong \Zp^{2n}$, $\Cliff/\Pauli \cong \PCliff/\PPauli \cong \Sp$, and $\PCliff \cong
\PPauli \rtimes \Sp$.  These ``glue facts'' are the two not formalised:
our semantic groups are \emph{defined} as the right-hand sides.
In~\cite{bian2026qudit} the global phases also include $-1$, so that
the swap and the multiplier gates belong to $\Cliff$; we work mainly
in the projective quotient, where $-1$ makes no difference except that
the relations come out a little simpler.

Using the framework of Appendix~\ref{sec:permutations}, we encode
circuits, circuit relations, and reasoning with the relations.  We
draw circuits in matrix-multiplication order, the reverse of the usual
circuit order: $w \bullet v$ applies $v$ first but draws it on the
right.  For example, the swap:

\noindent\begin{minipage}[c]{0.56\textwidth}
%% Generated by make agda from lagda/ex-circuit.lagda.tex (agda --latex --only-scope-checking);
%% edit that file, not this one.  The hidden block renders nothing.
\begin{code}[hide]%
\>[0]\AgdaKeyword{module}\AgdaSpace{}%
\AgdaModule{vqupit-short.lagda.ex-circuit}\AgdaSpace{}%
\AgdaKeyword{where}\<%
\\
%
\\[\AgdaEmptyExtraSkip]%
\>[0]\AgdaKeyword{open}\AgdaSpace{}%
\AgdaKeyword{import}\AgdaSpace{}%
\AgdaModule{vqupit-short.lagda.Prelude}\<%
\\
%
\\[\AgdaEmptyExtraSkip]%
\>[0]\AgdaKeyword{private}\AgdaSpace{}%
\AgdaKeyword{variable}\<%
\\
\>[0][@{}l@{\AgdaIndent{0}}]%
\>[2]\AgdaGeneralizable{n}\AgdaSpace{}%
\AgdaSymbol{:}\AgdaSpace{}%
\AgdaDatatype{ℕ}\<%
\end{code}
\begin{code}%
\>[0]\AgdaFunction{Ex}\AgdaSpace{}%
\AgdaSymbol{:}\AgdaSpace{}%
\AgdaDatatype{Word}\AgdaSpace{}%
\AgdaSymbol{(}\AgdaPostulate{Gen}\AgdaSpace{}%
\AgdaSymbol{(}\AgdaInductiveConstructor{₂₊}\AgdaSpace{}%
\AgdaGeneralizable{n}\AgdaSymbol{))}\<%
\\
\>[0]\AgdaFunction{Ex}\AgdaSpace{}%
\AgdaSymbol{=}\AgdaSpace{}%
\AgdaPostulate{CZ}\AgdaSpace{}%
\AgdaOperator{\AgdaInductiveConstructor{•}}\AgdaSpace{}%
\AgdaPostulate{H}\AgdaSpace{}%
\AgdaOperator{\AgdaFunction{↓}}\AgdaSpace{}%
\AgdaOperator{\AgdaInductiveConstructor{•}}\AgdaSpace{}%
\AgdaPostulate{H}\AgdaSpace{}%
\AgdaOperator{\AgdaPostulate{↑}}\AgdaSpace{}%
\AgdaOperator{\AgdaInductiveConstructor{•}}\AgdaSpace{}%
\AgdaPostulate{CZ}\AgdaSpace{}%
\AgdaOperator{\AgdaInductiveConstructor{•}}\AgdaSpace{}%
\AgdaPostulate{H}\AgdaSpace{}%
\AgdaOperator{\AgdaFunction{↓}}\AgdaSpace{}%
\AgdaOperator{\AgdaInductiveConstructor{•}}\AgdaSpace{}%
\AgdaPostulate{H}\AgdaSpace{}%
\AgdaOperator{\AgdaPostulate{↑}}\AgdaSpace{}%
\AgdaOperator{\AgdaInductiveConstructor{•}}\AgdaSpace{}%
\AgdaPostulate{CZ}\AgdaSpace{}%
\AgdaOperator{\AgdaInductiveConstructor{•}}\AgdaSpace{}%
\AgdaPostulate{H}\AgdaSpace{}%
\AgdaOperator{\AgdaFunction{↓}}\AgdaSpace{}%
\AgdaOperator{\AgdaInductiveConstructor{•}}\AgdaSpace{}%
\AgdaPostulate{H}\AgdaSpace{}%
\AgdaOperator{\AgdaPostulate{↑}}\<%
\end{code}

\end{minipage}\hfill
\begin{minipage}[c]{0.42\textwidth}\centering
\ufig[0.48]{def-Ex-eq}
\end{minipage}

\noindent A rule set is an inductive family indexed by the width, one
constructor per rule, each stated once at the smallest width at which
it fits; for example, two of the fifteen constructors of our primary
rule set \aid{{\_}SRel,{\_}==={\_}}, the order of the swap and the
rule moving $H$ across it:

\noindent\begin{minipage}[c]{0.56\textwidth}
%% Generated by make agda from lagda/ex-relation.lagda.tex (agda --latex --only-scope-checking);
%% edit that file, not this one.  The hidden block renders nothing.
\begin{code}[hide]%
\>[0]\AgdaKeyword{module}\AgdaSpace{}%
\AgdaModule{vqupit-short.lagda.ex-relation}\AgdaSpace{}%
\AgdaKeyword{where}\<%
\\
%
\\[\AgdaEmptyExtraSkip]%
\>[0]\AgdaKeyword{open}\AgdaSpace{}%
\AgdaKeyword{import}\AgdaSpace{}%
\AgdaModule{vqupit-short.lagda.Prelude}\<%
\\
\>[0]\AgdaKeyword{open}\AgdaSpace{}%
\AgdaKeyword{import}\AgdaSpace{}%
\AgdaModule{vqupit-short.lagda.ex-circuit}\<%
\\
%
\\[\AgdaEmptyExtraSkip]%
\>[0]\AgdaKeyword{private}\AgdaSpace{}%
\AgdaKeyword{variable}\<%
\\
\>[0][@{}l@{\AgdaIndent{0}}]%
\>[2]\AgdaGeneralizable{n}\AgdaSpace{}%
\AgdaSymbol{:}\AgdaSpace{}%
\AgdaDatatype{ℕ}\<%
\\
%
\\[\AgdaEmptyExtraSkip]%
\>[0]\AgdaKeyword{infix}\AgdaSpace{}%
\AgdaNumber{4}\AgdaSpace{}%
\AgdaOperator{\AgdaDatatype{\AgdaUnderscore{}CRel,\AgdaUnderscore{}===\AgdaUnderscore{}}}\<%
\end{code}
\begin{code}%
\>[0]\AgdaKeyword{data}\AgdaSpace{}%
\AgdaOperator{\AgdaDatatype{\AgdaUnderscore{}CRel,\AgdaUnderscore{}===\AgdaUnderscore{}}}\AgdaSpace{}%
\AgdaSymbol{:}\AgdaSpace{}%
\AgdaSymbol{(}\AgdaBound{n}\AgdaSpace{}%
\AgdaSymbol{:}\AgdaSpace{}%
\AgdaDatatype{ℕ}\AgdaSymbol{)}\AgdaSpace{}%
\AgdaSymbol{→}\AgdaSpace{}%
\AgdaFunction{WRel}\AgdaSpace{}%
\AgdaSymbol{(}\AgdaPostulate{Gen}\AgdaSpace{}%
\AgdaBound{n}\AgdaSymbol{)}\AgdaSpace{}%
\AgdaKeyword{where}\<%
\\
\>[0][@{}l@{\AgdaIndent{0}}]%
\>[2]\AgdaInductiveConstructor{order-Ex}%
\>[14]\AgdaSymbol{:}\AgdaSpace{}%
\AgdaSymbol{(}\AgdaInductiveConstructor{₂₊}\AgdaSpace{}%
\AgdaGeneralizable{n}\AgdaSymbol{)}\AgdaSpace{}%
\AgdaOperator{\AgdaDatatype{CRel,}}%
\>[30]\AgdaFunction{Ex}\AgdaSpace{}%
\AgdaOperator{\AgdaFunction{\textasciicircum{}}}\AgdaSpace{}%
\AgdaNumber{2}\AgdaSpace{}%
\AgdaOperator{\AgdaDatatype{===}}\AgdaSpace{}%
\AgdaInductiveConstructor{ε}\<%
\\
%
\>[2]\AgdaInductiveConstructor{semi-Ex-H↑}%
\>[14]\AgdaSymbol{:}\AgdaSpace{}%
\AgdaSymbol{(}\AgdaInductiveConstructor{₂₊}\AgdaSpace{}%
\AgdaGeneralizable{n}\AgdaSymbol{)}\AgdaSpace{}%
\AgdaOperator{\AgdaDatatype{CRel,}}%
\>[30]\AgdaFunction{Ex}\AgdaSpace{}%
\AgdaOperator{\AgdaInductiveConstructor{•}}\AgdaSpace{}%
\AgdaPostulate{H}\AgdaSpace{}%
\AgdaOperator{\AgdaPostulate{↑}}\AgdaSpace{}%
\AgdaOperator{\AgdaDatatype{===}}\AgdaSpace{}%
\AgdaPostulate{H}\AgdaSpace{}%
\AgdaOperator{\AgdaFunction{↓}}\AgdaSpace{}%
\AgdaOperator{\AgdaInductiveConstructor{•}}\AgdaSpace{}%
\AgdaFunction{Ex}\<%
\end{code}

\end{minipage}\hfill
\begin{minipage}[c]{0.42\textwidth}\centering
\ufig[0.48]{pv0-order-Ex}\\[3pt]
\ufig[0.48]{pv0-semi-Ex-Hup}
\end{minipage}

\noindent And reasoning with the rules is equational reasoning in the
congruence $\approx$ they generate, with combinators for symmetry,
associativity, and rewriting the left or the right factor of a
product; for example, conjugating by the swap moves $H$ down a wire:

\noindent\begin{minipage}[c]{0.56\textwidth}
%% Generated by make agda from lagda/ex-reasoning.lagda.tex (agda --latex --only-scope-checking);
%% edit that file, not this one.  The hidden block renders nothing.
\begin{code}[hide]%
\>[0]\AgdaKeyword{module}\AgdaSpace{}%
\AgdaModule{vqupit-short.lagda.ex-reasoning}\AgdaSpace{}%
\AgdaKeyword{where}\<%
\\
%
\\[\AgdaEmptyExtraSkip]%
\>[0]\AgdaKeyword{open}\AgdaSpace{}%
\AgdaKeyword{import}\AgdaSpace{}%
\AgdaModule{vqupit-short.lagda.Prelude}\<%
\\
\>[0]\AgdaKeyword{open}\AgdaSpace{}%
\AgdaKeyword{import}\AgdaSpace{}%
\AgdaModule{vqupit-short.lagda.ex-circuit}\<%
\\
\>[0]\AgdaKeyword{open}\AgdaSpace{}%
\AgdaKeyword{import}\AgdaSpace{}%
\AgdaModule{vqupit-short.lagda.ex-relation}\<%
\\
%
\\[\AgdaEmptyExtraSkip]%
\>[0]\AgdaKeyword{private}\AgdaSpace{}%
\AgdaKeyword{variable}\<%
\\
\>[0][@{}l@{\AgdaIndent{0}}]%
\>[2]\AgdaGeneralizable{n}\AgdaSpace{}%
\AgdaSymbol{:}\AgdaSpace{}%
\AgdaDatatype{ℕ}\<%
\\
%
\\[\AgdaEmptyExtraSkip]%
\>[0]\AgdaKeyword{postulate}\<%
\\
\>[0][@{}l@{\AgdaIndent{0}}]%
\>[2]\AgdaPostulate{axiom}\AgdaSpace{}%
\AgdaSymbol{:}\AgdaSpace{}%
\AgdaSymbol{∀}\AgdaSpace{}%
\AgdaSymbol{\{}\AgdaBound{w}\AgdaSpace{}%
\AgdaBound{v}\AgdaSpace{}%
\AgdaSymbol{:}\AgdaSpace{}%
\AgdaDatatype{Word}\AgdaSpace{}%
\AgdaSymbol{(}\AgdaPostulate{Gen}\AgdaSpace{}%
\AgdaGeneralizable{n}\AgdaSymbol{)\}}\AgdaSpace{}%
\AgdaSymbol{→}\AgdaSpace{}%
\AgdaGeneralizable{n}\AgdaSpace{}%
\AgdaOperator{\AgdaDatatype{SRel,}}\AgdaSpace{}%
\AgdaBound{w}\AgdaSpace{}%
\AgdaOperator{\AgdaDatatype{===}}\AgdaSpace{}%
\AgdaBound{v}\AgdaSpace{}%
\AgdaSymbol{→}\AgdaSpace{}%
\AgdaBound{w}\AgdaSpace{}%
\AgdaOperator{\AgdaPostulate{≈}}\AgdaSpace{}%
\AgdaBound{v}\<%
\end{code}
\begin{code}%
\>[0]\AgdaFunction{conj-Ex-H↑}\AgdaSpace{}%
\AgdaSymbol{:}\AgdaSpace{}%
\AgdaFunction{Ex}\AgdaSpace{}%
\AgdaOperator{\AgdaInductiveConstructor{•}}\AgdaSpace{}%
\AgdaPostulate{H}\AgdaSpace{}%
\AgdaOperator{\AgdaPostulate{↑}}\AgdaSpace{}%
\AgdaOperator{\AgdaInductiveConstructor{•}}\AgdaSpace{}%
\AgdaFunction{Ex}\AgdaSpace{}%
\AgdaOperator{\AgdaPostulate{≈}}\AgdaSpace{}%
\AgdaPostulate{H}\AgdaSpace{}%
\AgdaOperator{\AgdaFunction{↓}}\<%
\\
\>[0]\AgdaFunction{conj-Ex-H↑}\AgdaSpace{}%
\AgdaSymbol{=}\AgdaSpace{}%
\AgdaOperator{\AgdaPostulate{begin}}\<%
\\
\>[0][@{}l@{\AgdaIndent{0}}]%
\>[2]\AgdaFunction{Ex}\AgdaSpace{}%
\AgdaOperator{\AgdaInductiveConstructor{•}}\AgdaSpace{}%
\AgdaPostulate{H}\AgdaSpace{}%
\AgdaOperator{\AgdaPostulate{↑}}\AgdaSpace{}%
\AgdaOperator{\AgdaInductiveConstructor{•}}\AgdaSpace{}%
\AgdaFunction{Ex}%
\>[19]\AgdaOperator{\AgdaPostulate{≈⟨}}\AgdaSpace{}%
\AgdaPostulate{sym}\AgdaSpace{}%
\AgdaPostulate{assoc}\AgdaSpace{}%
\AgdaOperator{\AgdaPostulate{⟩}}\<%
\\
%
\>[2]\AgdaSymbol{(}\AgdaFunction{Ex}\AgdaSpace{}%
\AgdaOperator{\AgdaInductiveConstructor{•}}\AgdaSpace{}%
\AgdaPostulate{H}\AgdaSpace{}%
\AgdaOperator{\AgdaPostulate{↑}}\AgdaSymbol{)}\AgdaSpace{}%
\AgdaOperator{\AgdaInductiveConstructor{•}}\AgdaSpace{}%
\AgdaFunction{Ex}%
\>[19]\AgdaOperator{\AgdaPostulate{≈⟨}}\AgdaSpace{}%
\AgdaOperator{\AgdaPostulate{cleft}}\AgdaSpace{}%
\AgdaPostulate{axiom}\AgdaSpace{}%
\AgdaInductiveConstructor{semi-Ex-H↑}\AgdaSpace{}%
\AgdaOperator{\AgdaPostulate{⟩}}\<%
\\
%
\>[2]\AgdaSymbol{(}\AgdaPostulate{H}\AgdaSpace{}%
\AgdaOperator{\AgdaFunction{↓}}\AgdaSpace{}%
\AgdaOperator{\AgdaInductiveConstructor{•}}\AgdaSpace{}%
\AgdaFunction{Ex}\AgdaSymbol{)}\AgdaSpace{}%
\AgdaOperator{\AgdaInductiveConstructor{•}}\AgdaSpace{}%
\AgdaFunction{Ex}%
\>[19]\AgdaOperator{\AgdaPostulate{≈⟨}}\AgdaSpace{}%
\AgdaPostulate{assoc}\AgdaSpace{}%
\AgdaOperator{\AgdaPostulate{⟩}}\<%
\\
%
\>[2]\AgdaPostulate{H}\AgdaSpace{}%
\AgdaOperator{\AgdaFunction{↓}}\AgdaSpace{}%
\AgdaOperator{\AgdaInductiveConstructor{•}}\AgdaSpace{}%
\AgdaFunction{Ex}\AgdaSpace{}%
\AgdaOperator{\AgdaInductiveConstructor{•}}\AgdaSpace{}%
\AgdaFunction{Ex}%
\>[19]\AgdaOperator{\AgdaPostulate{≈⟨}}\AgdaSpace{}%
\AgdaOperator{\AgdaPostulate{cright}}\AgdaSpace{}%
\AgdaPostulate{axiom}\AgdaSpace{}%
\AgdaInductiveConstructor{order-Ex}\AgdaSpace{}%
\AgdaOperator{\AgdaPostulate{⟩}}\<%
\\
%
\>[2]\AgdaPostulate{H}\AgdaSpace{}%
\AgdaOperator{\AgdaFunction{↓}}\AgdaSpace{}%
\AgdaOperator{\AgdaInductiveConstructor{•}}\AgdaSpace{}%
\AgdaInductiveConstructor{ε}%
\>[19]\AgdaOperator{\AgdaPostulate{≈⟨}}\AgdaSpace{}%
\AgdaPostulate{right-unit}\AgdaSpace{}%
\AgdaOperator{\AgdaPostulate{⟩}}\<%
\\
%
\>[2]\AgdaPostulate{H}\AgdaSpace{}%
\AgdaOperator{\AgdaFunction{↓}}\AgdaSpace{}%
\AgdaOperator{\AgdaPostulate{∎}}\<%
\end{code}

\end{minipage}\hfill
\begin{minipage}[c]{0.42\textwidth}\centering
\ufig[0.48]{pv0-conj-Ex-Hup-steps}
\end{minipage}

\noindent Further, we can express and prove our main theorems:

\begin{theorem}[\aid{clifford-presentation}]\label{thm:main}
For every width $n$, the fifteen rules---the relations
of~\cite{bian2026qudit}, read modulo scalars---present the projective
Clifford group:
%% Copied from ../vqupit/latex/vqupit/sections/composing.tex (composing, the clifford-presentation block), the LaTeX
%% that `agda --latex --only-scope-checking' generated for the long version.
\begin{code}%
\>[0]\AgdaFunction{clifford-presentation}\AgdaSpace{}%
\AgdaSymbol{:}\AgdaSpace{}%
\AgdaSymbol{∀}\AgdaSpace{}%
\AgdaBound{n}\AgdaSpace{}%
\AgdaSymbol{→}\AgdaSpace{}%
\AgdaSymbol{(}\AgdaOperator{\AgdaPostulate{PapV1.Clifford-Relations.\AgdaUnderscore{}QRel,\AgdaUnderscore{}===\AgdaUnderscore{}}}\AgdaSpace{}%
\AgdaBound{n}\AgdaSymbol{)}\AgdaSpace{}%
\AgdaOperator{\AgdaRecord{IsPresentationOf}}\AgdaSpace{}%
\AgdaSymbol{(}\AgdaPostulate{Pauli⋊Sp}\AgdaSpace{}%
\AgdaBound{n}\AgdaSymbol{)}\<%
\end{code}

\end{theorem}

\noindent \aid{{\_}QRel,{\_}==={\_}} closes \aid{{\_}SRel,{\_}==={\_}}
under the three structural rules (Appendix~\ref{sec:permutations});
$\aid{Pauli⋊Sp}\;n$ is the semidirect product of $\Zp^{2n}$
by $\Sp$ (\S\ref{sec:composing}); \aid{{\_}IsPresentationOf{\_}}
(Appendix~\ref{sec:permutations}) is a group isomorphism $\mathsf{Cir}\;n/{\approx}
\to \PPauli \rtimes \Sp$:
\emph{soundness} (well-definedness of the map), each rule holds;
\emph{completeness}, $\sem{\_}$ is injective; and \emph{surjectivity},
each projective Clifford operator is represented by a circuit.

We refine the qupit paper's normal form for symplectic operators
(Clifford operators modulo Paulis) into a doubly inductive definition
(\S\ref{sec:symplectic}) and prove that it
enjoys the unique-normal-form property:

\begin{theorem}[\aid{unique-nf}]\label{thm:unique}
Normal forms with equal symplectic denotations are equal:
%% Copied from ../vqupit/latex/vqupit/sections/background.tex (background, block 4), the LaTeX
%% that `agda --latex --only-scope-checking' generated for the long version.
\begin{code}%
\>[0]\AgdaFunction{Theorem-2}\AgdaSpace{}%
\AgdaSymbol{:}\AgdaSpace{}%
\AgdaSymbol{∀}\AgdaSpace{}%
\AgdaBound{n}\AgdaSpace{}%
\AgdaSymbol{\{}\AgdaBound{u}\AgdaSpace{}%
\AgdaBound{v}\AgdaSpace{}%
\AgdaSymbol{:}\AgdaSpace{}%
\AgdaPostulate{NF}\AgdaSpace{}%
\AgdaBound{n}\AgdaSymbol{\}}\AgdaSpace{}%
\AgdaSymbol{→}\AgdaSpace{}%
\AgdaOperator{\AgdaPostulate{⟦}}\AgdaSpace{}%
\AgdaOperator{\AgdaPostulate{[}}\AgdaSpace{}%
\AgdaBound{u}\AgdaSpace{}%
\AgdaOperator{\AgdaPostulate{]}}\AgdaSpace{}%
\AgdaOperator{\AgdaPostulate{⟧}}\AgdaSpace{}%
\AgdaOperator{\AgdaPostulate{≈ˢ}}\AgdaSpace{}%
\AgdaOperator{\AgdaPostulate{⟦}}\AgdaSpace{}%
\AgdaOperator{\AgdaPostulate{[}}\AgdaSpace{}%
\AgdaBound{v}\AgdaSpace{}%
\AgdaOperator{\AgdaPostulate{]}}\AgdaSpace{}%
\AgdaOperator{\AgdaPostulate{⟧}}\AgdaSpace{}%
\AgdaSymbol{→}\AgdaSpace{}%
\AgdaBound{u}\AgdaSpace{}%
\AgdaOperator{\AgdaDatatype{≡}}\AgdaSpace{}%
\AgdaBound{v}\<%
\\
\>[0]\AgdaFunction{Theorem-2}\AgdaSpace{}%
\AgdaSymbol{=}\AgdaSpace{}%
\AgdaPostulate{U.⟦[]⟧-injective}\<%
\end{code}

\end{theorem}

\noindent\emph{The route.}  We mostly follow the paper: its normal
boxes, its symplectic normal form and its
parametrised \emph{box relations} --- pushing a gate through a normal
box yields an updated box and residual \emph{dirty} gates --- reappear
below, with Theorem~4.4's reduction to eighteen relations, a step the
paper already delegated to Agda.  The one major divergence is the step
that promotes symplectic to projective completeness (its
Proposition~4.9): the paper corrects each rule $C_l = C_r$ to $C_l =
C_r\,;P$ for the unique Pauli $P$ making it projectively sound and
replays the symplectic rewrites, pushing each correction to the end,
where it becomes ``a potentially different Pauli''.  We did not see
how to verify it at reasonable cost.  We instead work in a syntax where Paulis are
genuinely trivial --- symplectic circuits have only $H$, $S$, $CZ$ and
denote $\Sp$ itself --- and reintroduce Paulis and scalars by
construction (Figure~\ref{fig:route} in the appendix): the Pauli
bookkeeping of Proposition~4.9 is done once, by a semidirect-product
theorem proved for arbitrary groups before this project began.
